\chapter{DAGonWeb}
\label{app:dagonweb}

Every workflow in this book is a Python program. DAGonWeb
(\url{https://github.com/DagOnStar/DAGonWeb}) offers a second way in: a Flask
and Bootstrap Progressive Web Application in which a workflow is drawn on a
canvas, saved as a DAGonStar document and executed on the server by the
DAGonStar runtime. This appendix accompanies course deck~A4. It describes
DAGonWeb as documented in its repository (\texttt{README.md} and
\texttt{docs/}), and its example runs a DAGonWeb document with DAGonStar alone.

\begin{cautionbox}[Versions]
DAGonWeb pins DAGonStar~3.0.0 (commit \texttt{f62ceea}) in its
\texttt{requirements.txt}. Its palette offers all sixteen task types of that
version, including the volume tasks and meta-tasks of
Appendices~\ref{app:volumes} and~\ref{app:meta-tasks}. Check the pins of both
projects before a deployment.
\end{cautionbox}

\section{What DAGonWeb Does}

A user picks a template or starts from an empty canvas, adds tasks from a
palette, configures each one in an inspector and draws links between them.
Every change is saved automatically. \emph{Run} starts the workflow on the
server; the run page shows the graph updating live, the stage of each task
(staging inputs, executing, staging outputs), the captured error of a failed
task, and a file browser over each task's folder with previews for text, JSON,
CSV, GeoJSON and NetCDF. The palette and the inspector forms are generated from
one registry, \texttt{app/task\_specs.py}, which also builds the runtime
\texttt{DagonTask} of each task.

The category \emph{Volumes \& meta-workflows} holds the two task types new in
DAGonStar~3.0. A \emph{Volume} task creates an ad-hoc volume with the
\texttt{dagonfs}, \texttt{hercules}, \texttt{capio} or \texttt{dynostore}
provider and stages data into it, or destroys a volume and stages its results
out; the document stores the specification under the task's \texttt{volume}
key. A \emph{Meta} task runs a command-line, Nextflow, Snakemake, Parsl,
StreamFlow or CWL workflow unmodified; its specification is stored under the
\texttt{meta} key. Both are the keys that DAGonStar's \texttt{load\_json} reads,
and under portable emulation a local directory and the mock engine stand in
for the volume and the engine.

Five task types belong to DAGonWeb itself. \emph{Upload} and \emph{Form} pause
a run until a user uploads files or answers a form on the run page;
\emph{Table}, \emph{Chart} and \emph{Download} turn results into a table, a
chart or download buttons. In the document they are plain DAGonStar
\texttt{batch} tasks, whose command is a self-contained shell and Python
script, so the document still loads with DAGonStar.

\section{The Document}

\begin{definitionbox}{DAGonWeb document}
The DAGonStar task-map document, unchanged, plus additive \texttt{dagonweb}
keys. A task's \texttt{dagonweb} key holds its canvas position and form
values; the top-level \texttt{dagonweb.runtime} holds \texttt{max\_threads},
\texttt{portable\_emulation} and \texttt{pre\_stage}. DAGonStar's
\texttt{load\_json} ignores the extra keys.
\end{definitionbox}

Edges are always rebuilt from the document; stored \texttt{prevs} and
\texttt{nexts} are ignored. A \emph{data link}, drawn solid, is a
\texttt{workflow:///producer/file} reference: DAGonStar infers the edge and
stages the file (Lecture~08). An \emph{explicit dependency}, drawn dashed, is
listed in the consumer's \texttt{explicit\_prevs}: it orders the tasks without
moving data (Lecture~07). A named input of a batch task, such as
\texttt{obs = workflow:///fetch/outputs/observations.csv}, is prepended to the
command as a shell assignment, so the script uses \texttt{\$obs}.

\section{What Run Does}

When a user presses Run, DAGonWeb creates a run record and a folder
\texttt{run\_<id>/} below the scratch root, and dispatches the run to a thread
of the web process or, for many concurrent runs, to a Redis queue served by
worker processes. The run merges the global and the per-workflow
\texttt{dagon.ini}, forcing \texttt{scratch\_dir\_base} to the run folder;
builds a \texttt{Workflow} with the runtime options and a checkpoint file
\texttt{run\_<id>/\_dagonweb/checkpoint.json}; adds a \texttt{DagonTask} per
task; adds the explicit edges and calls \texttt{make\_dependencies()}; then
calls \texttt{launch()} and \texttt{wait()}. DAGonStar's event listeners write
each task's progress into the database. On a failed run, \emph{Resume} reuses
the same folder and checkpoint, so only the failed task and those after it run
again; \emph{Run again} starts afresh.

\Cref{lst:dagonweb-document} writes a DAGonWeb document for the station means
by hand and runs it the way DAGonWeb's executor does. The first run fails on a
typo in \texttt{report} (in DAGonStar~3.0 a non-zero exit status fails the task
and its dependants); the corrected document is then resumed from the
checkpoint.

\lstinputlisting[caption={A DAGonWeb document run by DAGonStar (\texttt{examples/A4-dagonweb/\allowbreak dagonweb\_document.py}).},label={lst:dagonweb-document}]{../../examples/A4-dagonweb/dagonweb_document.py}

\lstinputlisting[style=bash,caption={Recorded output of \texttt{dagonweb\_document.py}.}]{../../examples/A4-dagonweb/output/dagonweb_document.txt}

The resumed run reused \texttt{fetch} and the three means, and ran only
\texttt{report} and \texttt{notify}. The explicit edge from \texttt{report} to
\texttt{notify} survives \texttt{as\_json()}.

\section{Beyond the Canvas}

\paragraph{Exports and imports.} A workflow can be exported as a headless
application ZIP (a commented DAGonStar program, the JSON document, CWL, a
\texttt{dagon.ini} with secrets removed, \texttt{requirements.txt} and
\texttt{run.sh}), as a single Python program that needs only DAGonStar and
accepts \texttt{--resume}, \texttt{--input} and \texttt{--form}, as DAGonWeb or
portable DAGonStar JSON, or as CWL~v1.2 through \texttt{saveAsCWL}. DAGonWeb
imports JSON documents, ZIP applications and Python programs; Python is read by
static analysis and never executed.

\paragraph{Assistant and dry run.} When an administrator enables an
OpenAI-compatible LLM endpoint, an assistant with tools designs and edits a
draft of the workflow from a plain-language description; the canvas changes
only when the user applies the draft. A dry run checks the graph and the
references, runs \texttt{bash -n} on shell commands, compiles Python and builds
the DAGonStar tasks under portable emulation, without executing anything.
Visitors without an account can dry-run templates in a guest playground.

\paragraph{Deployment.} DAGonWeb runs with Docker Compose, with SQLite,
PostgreSQL or MySQL. Installed on the login node of a cluster, it runs every
workflow as the Unix account of the user who starts it: the user adds a key to
\texttt{\~{}/.ssh/authorized\_keys} that is restricted to the
\texttt{dagonweb-runner} program, Slurm jobs are submitted as that user, and
nothing executes inside the DAGonWeb container. Accounts carry roles and
permissions, sign-in can use single sign-on, every file path is checked
against the run folder, and the interface is built for WCAG~2.1~AA.

\begin{practicebox}
A DAGonWeb document is a DAGonStar document. A workflow drawn in the browser
can therefore be versioned, reviewed, run in continuous integration and
resumed on a cluster with the same tools used for the Python workflows of
this book.
\end{practicebox}
